% QMB_n02_zustand_messung.tex — Kapitel 2: Zustände, Messung, Zeit
% Quellen: Dok. 230, 307, 322, 334
\chapter{Zustände, Messung und Zeit}
\label{ch:zustand}

\section{Zustände als Modenkonfigurationen}

In der Standard-QM ist ein Zustand ein Vektor $|\psi\rangle$ in einem
Hilbertraum $\mathcal{H}$. In der FFGFT ist ein Zustand eine
Modenkonfiguration $(z, r, \theta)$ auf dem Träger $T^4/\mathbb{Z}_3$.
Die Verbindung ist bijektiv: Der Spektraloperator
\begin{equation}
  \hat F = \sum_k f_k\, |k\rangle\langle k|,
  \label{eq:spektral}
\end{equation}
mit Moden $k$ aus dem Gitter $D_4$ ist selbstadjungiert auf
$\mathcal{H} = L^2(T^4)\otimes\mathbb{C}^3$.\status{B} (Dok.~327)
Jede QM-Rechnung ist in FFGFT-Sprache abbildbar und umgekehrt;
die Formalismen liefern dieselben statistischen Vorhersagen.\status{B}

\subsection{Die Übersetzungstabelle}

\begin{center}
\small\renewcommand{\arraystretch}{1.25}
\begin{tabular}{p{3.4cm}p{4.2cm}}
\toprule
Hilbertraum-QM & FFGFT \\
\midrule
Qubit $|\psi\rangle = \alpha|0\rangle+\beta|1\rangle$ & Mode $(z,r,\theta)\in S^1\times\mathbb{R}_{>0}\times S^1$ \\
Bloch-Sphäre $S^2$ & Zylindrische Faser $(r,\theta)$ \\
Superposition & Intermediäre Torsion einer Mode \\
Tensorprodukt $|\psi\rangle\otimes|\phi\rangle$ & Unabhängige Modenkonfigurationen \\
Verschränkung & Topologische Verbindung (nicht reduzierbar) \\
Unitäre Evolution & Phasendrehung auf $T^4$ \\
Messung / Kollaps & Polübergang zu $z=\pm1$ \\
Born-Regel $P=|\alpha|^2$ & Statistisches Mittel über unbekannte Phase $\theta$ \\
\bottomrule
\end{tabular}
\end{center}

\section{Messung ohne Kollaps-Postulat}

Der Kollaps der Wellenfunktion ist in der Standard-QM ein Postulat
ohne mechanische Erklärung. In der FFGFT ist Messung eine
\emph{geometrische Wechselwirkung}: Das Messgerät treibt die Mode
auf das nächstgelegene stabile Extrem ($z=+1$ oder $z=-1$). Das ist
eine deterministische Bewegung im Modenraum — aber wegen der
unbekannten Anfangsphase $\theta$ statistisch ununterscheidbar
von Born-Regel-Verhalten:\status{B}
\begin{equation}
  P_{\mathrm{FFGFT}}(z=+1) = \frac{1+z}{2} = |\alpha|^2 = P_{\mathrm{QM}}(0).
\end{equation}
Die beiden Formalismen liefern identische statistische Vorhersagen
(Dok.~230). Interpretativ unterscheiden sie sich: Was die QM als Zufall
beschreibt, ist in der FFGFT Unkenntnis der genauen Trägerkonfiguration.

\section{Die Abwesenheit des Zeitoperators}

Paulis Einwand zeigt, dass es zu einem nach unten beschränkten Hamilton-
operator keinen selbstadjungierten, kanonisch konjugierten Zeitoperator
geben kann.\status{E} Die Zeit fehlt im Hilbertraum — nicht als Versehen,
sondern strukturell.

In der FFGFT ist das kein Mangel: $\tilde T$ ist kein Operator, sondern
eine modengebundene Eigenschaft des Trägers. Für Mode $i$ gilt
$\tilde T_i = 1/m_i$; verschiedene Moden tragen verschiedene Zeitskalen.
Superposition von Moden verschiedener Masse ergibt damit
\emph{inkompatible Zeitskalen} — eine strukturelle Aussage über die
Grenzen modenübergreifender Superposition.\status{B} (Dok.~334)

Die Schrödinger-Gleichung bleibt als effektive Beschreibung vollständig
gültig; die Zeit $t$ darin ist die globale Koordinate, die aus dem
Träger ausgerollt wird (Typ-I-Pullback, Dok.~333).

\begin{laierspur}
Dass die QM keinen Zeitoperator hat, klingt technisch — bedeutet aber:
Die Theorie kann über die Dauer von Prozessen nichts sagen, nur über
Zustände zu festen Zeitpunkten. Die FFGFT macht das explizit: Jedes
Teilchen hat eine eigene innere Uhr ($\tilde T = 1/m$). Zwei Teilchen
verschiedener Masse schwingen buchstäblich in verschiedenen Rhythmen.
\end{laierspur}
\quelle{Dok.~230, 307, 322, 327, 333, 334}
